\section{The moving kernel, its exact metric, and the original inverse bound}
\label{sec:real-pair-moving-quotient}
The local two-branch construction RLA1--35 now permits a single
finite-dimensional comparison across its three actual order strata.
The complete original operator and orthonormal frames are WCF1--6
\cite{WCF}; the certified point and four original labels remain
RPZ1--19 and RPE1--26 \cite{RPZ005,RPE005}.  All directional
exponents used below are the proved RPD23--30 \cite{RPD006}.
We derive each new matrix and metric identity.  The construction
contains an explicitly stated orthogonal projection; it does not
identify a projected extension with the original operator.

Fix an integer $D\ge1$, put $N=D+1$, and keep $s\in\{1,a_{\rm amp}\}$.
All neighborhoods below may depend on these fixed choices.  They are
not assertions uniform in the amplification or cutoff.
Keep $\rho_j$, the full mass, centers and physical frames of RLA21.
Let $\mathsf P_D$ denote orthogonal projection onto the first $D+1$
target frames, and define
\[
 \mathscr R_p=\mathsf P_D T_A|_{\mathcal P_{D+2}},\qquad
 R(p)_{ij}=\begin{cases}g_{j-i}(p)\rho_j/\rho_i,&j\ge i,\\
 0,&j<i,\end{cases}
 \quad 0\le i\le D,\quad0\le j\le D+2.
 \tag{RQT1}
\]
This entry formula follows by coefficient comparison in WCF3,
with the unchanged factor $g=e^{-4\omega}E_A$ from RLA2.
Projection discards exactly target degrees $D+1,D+2$; every
coefficient of the retained entries is the original coefficient.
On the actual order-$v$ source $\mathcal P_{D+v}$, $0\le v\le2$,
the original operator already has degree at most $D$, so that
projection acts as the identity on its output.

Partition the source frames after degrees zero and one:
\[
 R=[L\ H],\quad L\in\mathbb C^{N\times2},\quad
 H\in\mathbb C^{N\times N},\quad
 H_{ij}=g_{j+2-i}\rho_{j+2}/\rho_i\quad(0\le i,j<N),
 \tag{RQT2}
\]
where $g_k=0$ for $k<0$.  At $p_*$ the coefficients $g_0,g_1$
vanish and $g_2=-b\ne0$.  Thus $L(p_*)=0$ and $H(p_*)$ is
upper triangular with diagonal $-b\rho_{i+2}/\rho_i$.
It is exactly the original order-two quotient matrix, including
all $g_3,\ldots,g_{D+2}$.  Its nonzero determinant gives a
neighborhood on which $H$ is invertible.  Define the exact matrices
\[
 \begin{gathered}Z=H^{-1}L,\qquad M=(I_N+ZZ^*)^{-1},\qquad
 W=(I_2+Z^*Z)^{-1},\\
 K=\ker R=\left\{\binom{l}{-Zl}:l\in\mathbb C^2\right\}.
 \end{gathered}\tag{RQT3}
\]
The last equality follows from $R(l,h)=H(Zl+h)$, and proves
surjectivity as well as the exact kernel.  In particular $K(p_*)$
is the original $\mathcal P_1$ inside $\mathcal P_{D+2}$.
Off the order-two locus it is the kernel of the projected extension;
it is not called the kernel of the full $T_A$.

\subsection{The quotient and both exact norm transports}
The map from a source vector to its quotient coordinate is
$w=h+Zl$.  Vectors with that same $w$ differ by $K$.
Minimizing the original squared norm
$\|l\|^2+\|w-Zl\|^2$ gives
\[
 l=Z^*Mw,\quad h=Mw,\quad
 \|[l,h]\|_{\mathcal P_{D+2}/K}^2=w^*Mw,\quad
 R^\dagger=\binom{Z^*}{I_N}MH^{-1}.
 \tag{RQT4}
\]
For a direct proof, the minimum solves
$(I_2+Z^*Z)l=Z^*w$.  The identity
$(I_2+Z^*Z)Z^*=Z^*(I_N+ZZ^*)$ gives the displayed $l$;
then $h=w-Zl=Mw$.  Subtracting this pair from any other
representative gives $(u,-Zu)$, which is orthogonal to the
minimizer.  Its squared norm is nonnegative and vanishes only
at zero.  This proves the minimum, its uniqueness, and the norm.
It also proves $RR^\dagger=I_N$ and that $R^\dagger y$ is
the unique preimage of $y$ orthogonal to $K$.

Let square roots be the positive Hermitian roots, and set
\[
 Q_+=\binom{Z^*}{I_N}M^{1/2},\qquad
 Q_K=\binom{I_2}{-Z}W^{1/2},\qquad
 \widetilde A=RQ_+=HM^{-1/2},\qquad
 \widetilde A^{-1}=M^{1/2}H^{-1}.
 \tag{RQT5}
\]
Multiplication gives $Q_+^*Q_+=I_N$, $Q_K^*Q_K=I_2$ and
$Q_K^*Q_+=0$.  Their columns consequently form an orthonormal
basis of the entire original source space.  In particular
$Q_+Q_+^*+Q_KQ_K^*=I_{N+2}$.  This proves that $Q_+$ realizes
the quotient isometrically, with precisely the metric RQT4.
It is a proved coordinate isometry between these two stated
spaces, not a rescaling of the original measure or operator.

Write $m(p)$ and $M_R(p)$ for the smallest and largest singular
values of $R$ (there are exactly $N$ positive ones).  Then
\[
 RR^*=H(I_N+ZZ^*)H^*=\widetilde A\widetilde A^*,\quad
 m\ge\|H^{-1}\|^{-1},\quad
 M_R\le\|H\|\sqrt{1+\|Z\|^2}.
 \tag{RQT6}
\]
Indeed the middle positive matrix is at least $I_N$, and its
largest eigenvalue is $1+\|Z\|^2$; apply the squared norm
identity to each vector.  Shrinking the neighborhood so that
$\|H-H(p_*)\|\le(2\|H(p_*)^{-1}\|)^{-1}$ gives
$m\ge(2\|H(p_*)^{-1}\|)^{-1}$.  This follows by expanding
$(I+H(p_*)^{-1}(H-H(p_*)))^{-1}$ as its convergent geometric
matrix series.  Thus the moving quotient inverse is bounded
on that neighborhood, with the original finite matrix constant.

\subsection{Every actual-order domain maps into that quotient}
Let $J_v:\mathbb C^N\to\mathbb C^{N+2}$ insert a vector in
source positions $v,\ldots,D+v$, for $v=0,1,2$.
On the stratum of actual order $v$, these are the original
orthonormal quotient representatives for
$\mathcal P_{D+v}/\mathcal P_{v-1}$, so define
\[
 A_v=RJ_v,\qquad T_v=Q_+^*J_v,\qquad E_v=T_v^*T_v
 =J_v^*(I-Q_KQ_K^*)J_v,\qquad A_v=\widetilde A T_v.
 \tag{RQT7}
\]
These are exact morphisms, not just determinant identities.
The proof uses $RQ_K=0$ and the complete orthogonal decomposition
following RQT5.  On its stated stratum, $A_v$ is the actual
RLA30 matrix, with diagonal $g_v\rho_{j+v}/\rho_j$.
Consequently $T_v$ is invertible there.  The map induced by
orthogonal projection from the actual-order source into
$\mathcal P_{D+2}/K$ has squared source metric $E_v$.
The original source metric itself remains $I_N$.

Its exact determinant is
\[
 \det E_v=\frac{|g_v|^{2N}}{\det(RR^*)}
       \prod_{j=0}^D\frac{\rho_{j+v}^2}{\rho_j^2},\qquad
 g_v=\frac{(-i)^v\mu_v}{v!}.
 \tag{RQT8}
\]
Take squared absolute determinants of RQT7 and use RQT6.
This also proves that every norm correction between the actual
quotient and the moving quotient is retained, not suppressed.
At $p_*$, $Z=0$ and $J_2=Q_+$; hence $T_2=I_N$, $E_2=I_N$,
and $\widetilde A=A_2=H$.  At order zero or one, RQT7--8
give the full correction instead of asserting those metrics equal.

\subsection{Exactly two angles control the old source}
Let $J_{\rm tail}$ insert in the last two positions $D+1,D+2$,
the orthogonal complement of the old source $J_0\mathbb C^N$.
Split $Z$ into its first $N-2$ and last two rows:
\[
 Z=\binom{Z_{\rm top}}{Z_{\rm bot}},\qquad
 V=J_{\rm tail}^*Q_K=-Z_{\rm bot}W^{1/2}.
 \tag{RQT9}
\]
The singular values of $T_0$ are exactly
\[
 \underbrace{1,\ldots,1}_{N-2},\quad d_1,\ d_2,
 \qquad 1\ge d_1\ge d_2\ge0,
 \qquad (d_1,d_2)=\operatorname{singularvalues}(V).
 \tag{RQT10}
\]
To prove it, put $F=J_0^*Q_K$.  Then $E_0=I_N-FF^*$.
The nonzero eigenvalues of $FF^*$ and $F^*F$ agree, with
multiplicity: the maps $F$ and $F^*$ intertwine their eigenspaces
at each nonzero eigenvalue and are inverse there up to that
eigenvalue.  Meanwhile
$F^*F=Q_K^*(I-J_{\rm tail}J_{\rm tail}^*)Q_K=I_2-V^*V$.
The remaining $N-2$ eigenvalues of $E_0$ are one.  This proves
RQT10 including endpoints and multiplicities, also when an extra
one or zero occurs.  Geometrically $d_j$ are the sines of the
two angles between $K$ and the old source: the formula above is
their exact definition through orthogonal projections in the
original Gamma metric.  Both are zero at $p_*$.

There is an exact two-by-two inverse calculation with all higher
coefficients present.  A vector in the old source has full source
coordinates $(l,m,0_2)$, so
\[
 H^{-1}B_{D,s}(l,m)=\binom{Z_{\rm top}l+m}{Z_{\rm bot}l},\qquad
 \det Z_{\rm bot}=\frac{g_0^N}{\det H}.
 \tag{RQT11}
\]
Here the determinant is valid also at singular parameters.
For proof, permute the old columns from $(l,m)$ to $(m,l)$;
the sign is $(-1)^{2(N-2)}=1$.  The resulting matrix is
$\bigl(\begin{smallmatrix}I&Z_{\rm top}\\0&Z_{\rm bot}\end{smallmatrix}\bigr)$,
and $\det B_{D,s}=g_0^N$.  At $g_0\ne0$, if
$H^{-1}y=(w_{\rm top},w_{\rm bot})$, its exact inverse is
\[
 l=Z_{\rm bot}^{-1}w_{\rm bot},\qquad
 m=w_{\rm top}-Z_{\rm top}Z_{\rm bot}^{-1}w_{\rm bot}.
 \tag{RQT12}
\]
Substitution proves both inverse compositions.  In particular
every divergence is accounted for by this explicit two-by-two
matrix and the bounded matrices around it; no coefficient has
been dropped.  The exact product of the two angle sines is
\[
 d_1d_2=\frac{|g_0|^N}{\sqrt{\det(RR^*)}}
 =\frac{|U(0,p)|^N|\alpha\beta|^N}
 {|\det H|\sqrt{\det(I_2+Z^*Z)}}.
 \tag{RQT13}
\]
This follows from RQT8 with $v=0$ and RQT10.  The second form
uses RLA7 and RQT6; the equality
$\det(I_N+ZZ^*)=\det(I_2+Z^*Z)$ follows from the same
nonzero-eigenvalue correspondence used in RQT10.

For $1\le r\le N$ put $\mathcal D_1=d_2^{-1}$ and
$\mathcal D_r=(d_1d_2)^{-1}$ for $r\ge2$.
Off the branches, the exact factorization and bounds are
\[
 B_{D,s}^{-1}=T_0^{-1}\widetilde A^{-1},\qquad
 \frac{\mathcal D_r}{\|\bigwedge^r\widetilde A\|}
 \le\|\bigwedge^r B_{D,s}^{-1}\|
 \le\mathcal D_r\|\bigwedge^r\widetilde A^{-1}\|.
 \tag{RQT14}
\]
The $r$th exterior norm of $T_0^{-1}$ is $\mathcal D_r$ by
RQT10, since exterior singular values are the products of the
corresponding ordinary singular values.  Submultiplicativity
applied to RQT7 gives the upper bound; applying it to
$T_0^{-1}=B_{D,s}^{-1}\widetilde A$ gives the lower one.
All multipliers here are finite original matrix norms.  They
may be bounded further by $m^{-r}$ and $M_R^r$ via RQT6.

\subsection{The literal relation to inverse-exterior equation (11)}
The inverse-exterior estimate under investigation follows from
the determinant and forward estimate WCF1--6 \cite{WCF}.
We restate its exact form here, retaining its actual stratum.
Put $b_s=s/2$, $n_v=D+v$ and
\[
\begin{gathered}
 A_{\rm abs}(p)=\sum_{u,w=0}^{8}|a_{uw}(p)|,\qquad
 R_{n_v,s}=\begin{cases}(4n_v+1)^{1/4},&s=1,\\1,&s=a_{\rm amp},\end{cases}\\
 U_{D,s}^{(v)}=e^{1+2\pi\gamma}A_{\rm abs}(p)
                 R_{n_v,s}(D+1)(2n_v+1)^{4\delta}.
\end{gathered}\tag{RQT15}
\]
For completeness the forward estimate is obtained on $|t|=r<1$
by bounding each original power
$((1-it)/(1+it))^{(u-4)\delta+i(w-4)\gamma}$ by
$e^{2\pi\gamma}((1+r)/(1-r))^{4\delta}$.
The fraction has positive real part, modulus between
$(1-r)/(1+r)$ and its reciprocal, and argument of absolute
value less than $\pi/2$.  Cauchy's coefficient formula and
$r=m/(m+1)$ give
$|g_m|\le e^{1+2\pi\gamma}A_{\rm abs}(2m+1)^{4\delta}$
for $m\ge1$; the constant coefficient obeys the same bound.
Use $m=j+v$ in the actual order matrix.  Its weight ratios
are at most $R_{n_v,s}$: for $s\ge2$ the sequence $\rho_j$
decreases; for $s=1$ it increases and
$\rho_n\le(4n+1)^{1/4}$ follows by induction from
$(4n+1)(n+1)^2\le(4n+5)(n+1/2)^2$.
Row and column sums are therefore at most $U_{D,s}^{(v)}$.
The weighted Cauchy--Schwarz inequality
$|\sum_j a_{ij}x_j|^2\le(\sum_j|a_{ij}|)
\sum_j|a_{ij}||x_j|^2$, summed in $i$, proves
$\|A_v\|\le U_{D,s}^{(v)}$.

Multiplying the singular values and the triangular diagonal gives
the requested bound, and RQT8 rewrites it exactly as
\[
\begin{aligned}
 \|\bigwedge^r A_v^{-1}\|
 &=\frac{\|\bigwedge^{N-r} A_v\|}{|\det A_v|}\\
 &\le (U_{D,s}^{(v)})^{N-r}
       \left|\frac{\mu_v}{v!}\right|^{-N}
       \prod_{j=0}^D\prod_{k=1}^v
         \left(\frac{j+b_s+k-1}{j+k}\right)^{1/2}\\
 &=\frac{(U_{D,s}^{(v)})^{N-r}}
       {\sqrt{\det(RR^*)}\sqrt{\det E_v}}.
\end{aligned}\tag{RQT16}
\]
For the first equality, if the forward singular values are
$\sigma_1\ge\cdots\ge\sigma_N>0$, both sides are
$(\sigma_{N-r+1}\cdots\sigma_N)^{-1}$.
For the last equality, take the positive square root of RQT8.
Thus the small divisor in equation (11) is exactly the product
of the moving-quotient volume and its section-volume correction.
On the old order-zero source, $\sqrt{\det E_0}=d_1d_2$.
At the certified order-two point, $E_2=I$ and
$\mu_2/2=b>0$, so equation (11) is finite for its actual domain.
There is no substitution of a zero $\mu_0$ into an order-two
formula.  Equally, the bounded moving quotient does not bound
the old inverse without its singular section correction.

On any real nonzero direction $v$ in the original $(z,x)$
coordinates, $g_0(p_*+\tau v)=c(v)\tau^2+O(\tau^3)$ with
$c(v)=|A_zv_z+A_xv_x|^2>0$.  Therefore RQT13 proves
\[
 \lim_{\tau\to0}|\tau|^{-2N}d_1d_2
 =\frac{c(v)^N}{b^N}\prod_{j=0}^D\frac{\rho_j}{\rho_{j+2}}.
 \tag{RQT17}
\]
The constant is the full inverse determinant of $H(p_*)$
times $c(v)^N$; its weights have not been omitted.
RQT14 and RPD23--30 also give the separate angle exponents:
away from a resonant direction they are $(D,D+2)$;
on a resonant straight direction they are $(D+1,D+1)$;
on the previously constructed turning paths with $0<q<1$
they are $(D+q,D+2-q)$, in descending angle order.
To justify transfer of each exponent, multiplication by the
invertible $\widetilde A$ and its bounded inverse bounds each
singular value by fixed positive factors.  This follows from
the minimum--maximum characterization of singular values applied
to the norm inequalities
$m\|T_0x\|\le\|B_{D,s}x\|\le M_R\|T_0x\|$.
RQT17 fixes their product constant; individual constants are
computed by the exact matrices RQT9--12, not asserted unchanged
by multiplication by $\widetilde A$.

\subsection{The four ES states and the scalar endpoint}
For $D=3$ let $X_0$ be the matrix of the four Lagrange polynomials
$\ell_a$ in the original source frames, ordered
$\mathrm M,\mathrm N,\mathrm T,\mathrm E$.  Their images in the
single moving quotient, and their original outputs, are exactly
\[
 \widehat X_0=T_0X_0,\qquad
 Y_0=B_{3,s}X_0=\widetilde A\widehat X_0.
 \tag{RQT18}
\]
No formal label changes.  Since $X_0$ is invertible, the only
rank loss here is the computed rank loss of $T_0$, which at
$p_*$ is two.  RLA31--35 recovers the two lost scalar
coordinates and their complete metric.  On each actual stratum,
the corrected inputs $[S^v\ell_a]$ have their own invertible
frame matrix $X_v$; exactly the same construction gives
$Y_v=\widetilde A T_vX_v$.  This proves the exact bridge from
all four ES states to RQT16, while preserving both the fixed
old inputs and the original order-adjusted inputs.

The case $D=0$ needs no artificial two-angle interpretation.
Here $R=(g_0,g_1\rho_1/\rho_0,g_2\rho_2/\rho_0)$ is one row,
its moving quotient has dimension one and
\[
 d_0=\frac{|g_0|}{\|R\|},\qquad
 |B_{0,s}^{-1}|=\frac1{|g_0|},\qquad
 \det E_v=\frac{|g_v|^2\rho_v^2}{\rho_0^2\|R\|^2}
 \quad(v\text{ the actual order}).
 \tag{RQT19}
\]
These are obtained by projecting each one-dimensional coordinate
section onto the span of $R^*$.  At $p_*$, $R$ has only its
last entry nonzero and $E_2=1$.  Thus the scalar endpoint and
every higher cutoff are included without changing their dimensions.
